\chapter*{Conclusioni}
\addcontentsline{toc}{chapter}{Conclusioni}

Questa tesina ha implementato e confrontato cinque approcci al Minimum Vertex Cover per rispondere alle tre domande poste nell'introduzione: forza bruta, branch-and-bound, programmazione lineare intera risolta con Gurobi, 2-approssimazione via matching massimale ed euristica greedy.

\begin{itemize}
    \item \textbf{Dove si trova il muro esponenziale.} La forza bruta si è rivelata impraticabile già nel pilot (Capitolo~4) ed è stata esclusa dalla campagna vera e propria. Tra i due metodi esatti rimasti, il muro pratico cade a dimensioni molto diverse: il branch-and-bound smette di essere un confronto utile già a $n\approx 40$-$80$, prima sui grafi densi che sugli sparsi, mentre l'ILP di Gurobi resta praticabile fino a $n\approx 200$-$300$ indipendentemente dalla densità. Il vantaggio dell'ILP non viene da una gestione più furba della densità del grafo, ma da un rilassamento LP che resta stringente in entrambi i regimi, a differenza del lower bound basato su matching massimale del branch-and-bound.
    \item \textbf{Andamento temporale.} Gli esperimenti hanno confermato l'analisi teorica. L'algoritmo forza bruta ha confermato l'andamento esponenziale. Il branch-and-bound ha confermato l'andamento esponenziale, ma con un esponente più piccolo, e con una forte dipendenza dalla densità del grafo. L'ILP di Gurobi ha confermato l'andamento esponenziale, ma con un esponente ancora più piccolo e senza dipendenza dalla densità. Le due euristiche hanno confermato l'andamento polinomiale, con il greedy leggermente più veloce del 2-approssimato.    
    \item \textbf{Qualità delle soluzioni.} Il 2-approssimato ha confermato empiricamente il proprio bound teorico su 540 istanze, senza mai superare un rapporto di 1.667 rispetto all'ottimo, ben al di sotto del limite di 2. Il greedy, pur senza alcuna garanzia teorica, si è però rivelato sistematicamente più vicino all'ottimo, con una differenza di velocità tra i due trascurabile rispetto al confronto con l'ILP.
\end{itemize}

